\section{Related Work}
\label{sec:related}

\subsection{Weight-Space Model Merging and Model Soups}
Weight-space model merging has gained prominent traction as a zero-shot multi-task consolidation strategy. Early approaches, such as Model Soups \cite{wortsman2022model} and Stochastic Weight Averaging (SWA) \cite{izmailov2018averaging}, demonstrate that averaging the weights of neural networks fine-tuned from the same pre-trained initialization on similar tasks can lead to flat minima and improved generalization. These techniques, however, assume that the fine-tuned models reside within the same local optimization basin—a phenomenon known as linear mode connectivity \cite{frankle2020linear}. When merging models fine-tuned on highly disparate or conflicting tasks, linear mode connectivity often breaks down, resulting in a severe performance drop.

To extend merging to disparate tasks, \citet{ilharco2022editing} introduced the concept of \emph{task vectors}, which are computed as the element-wise difference between fine-tuned and pre-trained weights ($V = W - W_{base}$). Task Arithmetic demonstrates that these task vectors can be added to, or subtracted from, the pre-trained base model to selectively steer model capabilities. Despite its simplicity and parameter efficiency, standard Task Arithmetic struggles under multi-task scenarios due to parameter interference and representational drift.

\subsection{Parameter Interference and Heuristic Resolution}
To address parameter interference, several state-of-the-art merging methods propose "denoising" task vectors prior to interpolation. TIES-Merging \cite{yadav2023ties} identifies that sign disagreement (where different experts update the same coordinate in opposite directions) and redundant parameter drift are the primary drivers of model degradation. TIES resolves this by pruning parameters of small magnitudes, determining a consensus sign across tasks for each parameter coordinate, and then averaging only those updates that agree with the consensus sign. Similarly, Sparse Task Arithmetic (STA) \cite{polymerge, zipmerge} applies a hard, layer-wise magnitude-based pruning boundary (e.g., keeping only the top $50\%$ largest values) before linear combination. 

While these heuristic methods are effective in practice, they operate strictly on a coordinate-wise basis. By treating each weight coordinate as independent, they disregard the low-rank correlation structures that naturally arise in overparameterized neural networks. Furthermore, coordinate-wise hard thresholding is non-differentiable and does not possess a closed-form geometric formulation, precluding any mathematical guarantees on the distortion of the representation spaces. In contrast, GSC-Merge replaces these heuristics with a mathematically rigorous, differentiable spectral consensus mechanism operating on the Grassmannian manifold.

\subsection{Tuning Merging Coefficients and Few-Shot Adaptation}
To optimize model performance post-merging, research has shifted from static, global scaling factors to learnable, layer-wise coefficient configurations. AdaMerging \cite{liang2023adamerging} proposes optimizing merging weights using test-time entropy minimization. However, fully test-time adaptation (TTA) methods like Tent \cite{wang2020tent} often suffer from task suite bias and transductive overfitting to stream noise \cite{liao2021are}. To prevent these limitations, OFS-Tune \cite{ofstune} and PolyMerge \cite{polymerge} introduce Offline Few-Shot Validation Tuning, wherein layer-wise merging coefficients are optimized on a tiny, labeled calibration set prior to evaluation.

Our work builds upon this few-shot tuning framework. We identify a major failure mode of unconstrained coefficient search on small validation sets, which we call the \emph{Overfitting-Optimizer Paradox}: without proper regularization, the learning process rapidly memorizes validation noise, causing the representation to drift away from the expert manifolds and degrading test performance. We demonstrate that projecting task updates onto a low-rank Grassmannian subspace provides an elegant, mathematically sound spectral regularizer that overcomes this paradox.

\subsection{Low-Rank Adaptations and Geometric Manifolds}
The geometry of low-rank subspaces has been widely exploited during training. Low-Rank Adaptation (LoRA) \cite{hu2021lora} constrains parameter updates during fine-tuning to a low-rank factorization, showing that downstream adaptation occurs in a very low intrinsic dimension. While LoRA enforces a low-rank constraint \emph{during training}, GSC-Merge is a post-hoc merging framework that identifies a shared low-rank \emph{consensus subspace} from fully fine-tuned, dense parameters. 

Our formulation connects model merging with the rich mathematics of the Grassmannian manifold—the set of all $r$-dimensional linear subspaces in a vector space \cite{edelman1998geometry}. By using Singular Value Decomposition (SVD) to construct the projection matrix, GSC-Merge performs an optimal low-rank reconstruction on the Grassmannian, ensuring that we preserve the maximum possible multi-task update energy while discarding destructive, incoherent noise.
